\documentclass[11pt,a4paper]{article}

% ----------------------------------------------------------------------
%  Philips Document Flow (PDF) - Manual Tecnico
%  Compila en Overleaf con pdfLaTeX (motor por defecto).
% ----------------------------------------------------------------------

\usepackage[utf8]{inputenc}
\usepackage[T1]{fontenc}
\usepackage[spanish,es-noshorthands]{babel}
\usepackage{lmodern}
\usepackage[a4paper,margin=2.5cm]{geometry}
\usepackage{xcolor}
\usepackage{booktabs}
\usepackage{array}
\usepackage{ragged2e}
\usepackage{tabularx}
\usepackage{xltabular}
\usepackage{enumitem}
\usepackage{listings}
\usepackage{fancyhdr}
\usepackage{titlesec}
\usepackage{tikz}
\usetikzlibrary{arrows.meta,positioning}
\usepackage[hidelinks]{hyperref}

% ---- Paleta sobria (texto siempre en negro) --------------------------
\definecolor{codebg}{RGB}{246,247,248}
\definecolor{linegray}{RGB}{200,205,210}
\definecolor{headgray}{RGB}{238,240,242}
\color{black}

% ---- Estilo de codigo (monocromo, texto negro) -----------------------
\lstdefinestyle{code}{
  basicstyle=\ttfamily\small\color{black},
  backgroundcolor=\color{codebg},
  frame=single,
  rulecolor=\color{linegray},
  breaklines=true,
  showstringspaces=false,
  columns=fullflexible,
  keepspaces=true,
  captionpos=b,
  xleftmargin=6pt,
  xrightmargin=6pt,
  aboveskip=10pt,
  belowskip=6pt,
}
\lstset{style=code}

% ---- Columnas de tabla que ajustan y NO se solapan -------------------
% L = columna flexible con texto alineado a la izquierda y ajuste de linea.
\newcolumntype{L}{>{\RaggedRight\arraybackslash}X}
\newcolumntype{C}{>{\Centering\arraybackslash}X}

% ---- Encabezado y pie ------------------------------------------------
\pagestyle{fancy}
\fancyhf{}
\setlength{\headheight}{14pt}
\fancyhead[L]{\small\itshape Philips Document Flow (PDF) --- Manual Técnico}
\fancyhead[R]{\small\thepage}
\renewcommand{\headrulewidth}{0.4pt}

% ---- Formato de titulos (negro, con linea fina) ----------------------
\titleformat{\section}
  {\Large\bfseries\color{black}}{\thesection}{0.6em}{}[\vspace{1pt}{\color{linegray}\titlerule}]
\titleformat{\subsection}{\large\bfseries\color{black}}{\thesubsection}{0.6em}{}
\titlespacing*{\section}{0pt}{16pt}{8pt}

% ---- Macros ----------------------------------------------------------
% Backslash con punto de corte de linea, para rutas largas en tablas.
\newcommand{\bs}{\textbackslash\allowbreak}
\newcommand{\ruta}[1]{\texttt{\small #1}}

\setlength{\parskip}{4pt}
\setlength{\parindent}{0pt}

\begin{document}

% ======================================================================
%  PORTADA
% ======================================================================
\begin{titlepage}
\centering
\vspace*{3cm}
{\color{black}\rule{\textwidth}{1.2pt}}\\[0.8cm]
{\Huge\bfseries Philips Document Flow (PDF)}\\[0.5cm]
{\LARGE Manual Técnico}\\[0.4cm]
{\color{black}\rule{\textwidth}{1.2pt}}\\[2cm]

{\large Aplicación de escritorio para la apertura automática de\\
documentos PDF descargados (instrucciones de trabajo WI / MPI)}\\[3cm]

\begin{tabular}{@{}r l@{}}
\textbf{Plataforma} & Microsoft Windows 10 / 11 (x64) \\[5pt]
\textbf{Tecnología} & .NET 8 --- Windows Forms (C\#) \\[5pt]
\textbf{Motor de visualización} & Microsoft WebView2 (Chromium) \\[5pt]
\textbf{Versión de la aplicación} & 2.2.0 \\[5pt]
\textbf{Fecha} & Junio de 2026 \\
\end{tabular}

\vfill
\end{titlepage}

% ======================================================================
\tableofcontents
\newpage

% ======================================================================
\section{Introducción}
% ======================================================================
PDF Auto Viewer es una aplicación de escritorio para Windows que automatiza la
visualización de documentos PDF descargados. Se ejecuta de forma silenciosa en
la bandeja del sistema y vigila la carpeta de descargas del usuario; cuando
aparece un documento nuevo lo abre en un visor integrado y, una vez que el
usuario termina de consultarlo, lo elimina automáticamente para mantener la
carpeta limpia.

La aplicación está orientada a entornos donde los documentos pueden distribuirse
en dos idiomas, identificados por un sufijo en el nombre del archivo:

\begin{itemize}[leftmargin=1.4em,itemsep=2pt]
  \item \texttt{\_SPA} --- documento en español.
  \item \texttt{\_ENG} --- documento en inglés.
\end{itemize}

El diseño es deliberadamente acotado: el sistema se limita a \textbf{vigilar},
\textbf{abrir}, \textbf{cerrar tras el uso} y \textbf{eliminar} los documentos.
No realiza ninguna otra acción sobre el equipo ni sobre los datos del usuario.

% ======================================================================
\section{Características principales}
% ======================================================================
\begin{itemize}[leftmargin=1.4em,itemsep=3pt]
  \item Apertura automática de los PDF descargados, sin intervención del usuario.
  \item Visor integrado propio basado en el motor WebView2 (Chromium).
  \item Eliminación automática del documento al cerrarse.
  \item Límite de visualización de 20 minutos, con un aviso previo a los 15.
  \item Reemplazo automático de un documento por su copia más reciente.
  \item Filtro de idioma configurable cuando se descargan ambas versiones.
  \item Prioridad por tipo: el PDF derivado de un \texttt{.docx}
        (\texttt{\_docx.pdf}) prevalece sobre el \texttt{.pdf} nativo del mismo
        documento e idioma.
  \item Detección robusta de descargas mediante eventos del sistema de archivos
        y una reexploración periódica de respaldo.
  \item Identidad visual propia: icono azul con el texto «PDF» aplicado al
        ejecutable, a la ventana, a la barra de tareas y a la bandeja del sistema.
  \item Interfaz de usuario en inglés, como estándar para su uso en distintas
        plantas.
\end{itemize}

% ======================================================================
\section{Arquitectura del sistema}
% ======================================================================
La aplicación se organiza en dos capas: una capa de lógica (\texttt{Core}) y una
capa de interfaz de usuario (\texttt{UI}). No existe una ventana principal
permanente; el punto de permanencia de la aplicación es el icono de la bandeja
del sistema.

\vspace{0.4cm}
\begin{center}
\begin{tikzpicture}[
  node distance=0.85cm and 1.25cm,
  every node/.style={font=\small},
  comp/.style={rectangle, rounded corners, draw=black, fill=white,
               text width=3.1cm, align=center, minimum height=1.0cm},
  ext/.style={rectangle, draw=black, dashed, fill=white, text width=2.7cm,
              align=center, minimum height=1.0cm},
  arr/.style={-{Stealth[length=2.4mm]}, thick}
]
\node[ext]  (fs)   {Carpeta de descargas};
\node[comp] (mon)  [right=of fs]  {FolderMonitor\\[2pt]\footnotesize vigilancia};
\node[comp] (life) [right=of mon] {PdfLifecycleManager\\[2pt]\footnotesize orquestación};
\node[comp] (view) [right=of life]{PdfViewerForm\\[2pt]\footnotesize visor (WebView2)};
\node[comp] (tray) [below=1.3cm of mon] {TrayApp\\[2pt]\footnotesize bandeja};
\node[comp] (set)  [left=of tray] {AppSettings\\[2pt]\footnotesize configuración};

\draw[arr] (fs)   -- (mon);
\draw[arr] (mon)  -- (life);
\draw[arr] (life) -- (view);
\draw[arr] (set)  -- (tray);
\draw[arr] (tray) -- (mon);
\draw[arr] (life) to[bend right=18] node[midway,below,font=\footnotesize]{eventos} (tray);
\end{tikzpicture}
\end{center}
\vspace{0.3cm}

\begin{xltabular}{\textwidth}{@{} >{\bfseries}p{5.1cm} L @{}}
\toprule
\normalfont\textbf{Componente} & \textbf{Responsabilidad} \\
\midrule
\endfirsthead
\toprule
\normalfont\textbf{Componente} & \textbf{Responsabilidad} \\
\midrule
\endhead
\bottomrule
\endfoot
Program & Punto de entrada. Inicializa el contexto de aplicación. \\
TrayApp & Mantiene viva la aplicación a través del icono de bandeja y conecta los componentes entre sí. \\
AppSettings & Carga y guarda la configuración del usuario en formato JSON. \\
FolderMonitor & Detecta la aparición de nuevos archivos PDF en la carpeta vigilada. \\
PdfLifecycleManager & Orquesta el ciclo de vida completo de cada documento detectado. \\
PdfViewerForm & Visor integrado basado en WebView2; muestra el PDF y controla su cierre. \\
StatusForm & Ventana principal: estado del monitor y selección de idioma. \\
\bottomrule
\end{xltabular}

% ======================================================================
\section{Flujo de procesamiento}
% ======================================================================
Cada documento detectado se procesa en una tarea en segundo plano independiente,
de modo que varias descargas simultáneas no se bloquean entre sí. El ciclo consta
de las siguientes etapas:

\begin{enumerate}[leftmargin=1.6em,itemsep=3pt]
  \item \textbf{Detección.} La carpeta vigilada notifica la aparición de un
        archivo con extensión \texttt{.pdf}.
  \item \textbf{Estabilización.} La aplicación espera a que el navegador termine
        de escribir el archivo, comprobando que el tamaño se mantenga estable y
        que ningún proceso conserve un bloqueo de escritura.
  \item \textbf{Apertura.} El documento se abre en el visor integrado.
  \item \textbf{Espera.} El ciclo se bloquea hasta que la ventana se cierra, ya
        sea por el usuario, por el reemplazo de una copia más reciente, por el
        filtro de idioma o por el límite de 20 minutos.
  \item \textbf{Eliminación.} El archivo se elimina de forma segura.
\end{enumerate}

\subsection{Control de duplicados y reprocesamiento}
La aplicación mantiene varias estructuras en memoria para garantizar que cada
documento se procese una sola vez y que no se pierda ninguna descarga:

\begin{itemize}[leftmargin=1.4em,itemsep=2pt]
  \item \textbf{En proceso:} rutas que se están procesando en ese momento.
  \item \textbf{Enfriamiento:} rutas finalizadas recientemente, para ignorar
        eventos tardíos del sistema de archivos.
  \item \textbf{Procesados:} rutas ya atendidas junto con su fecha de
        modificación, para no reabrir un documento que no ha cambiado.
\end{itemize}

% ======================================================================
\section{Detección de descargas}
% ======================================================================
La detección de archivos nuevos combina dos mecanismos complementarios que, en
conjunto, garantizan que ningún documento quede sin abrirse.

\subsection{Eventos del sistema de archivos}
El componente \texttt{FolderMonitor} utiliza la clase estándar
\texttt{FileSystemWatcher} de .NET, que notifica de forma inmediata la creación y
el renombrado de archivos. Los navegadores descargan a un archivo temporal (por
ejemplo \texttt{.crdownload}) y lo renombran a \texttt{.pdf} al completar la
descarga; ambos casos se detectan.

\subsection{Reexploración periódica de respaldo}
Como red de seguridad ante eventos perdidos (desbordamiento del búfer del sistema
o condiciones de carrera de los navegadores), un temporizador reexplora la
carpeta cada tres segundos y vuelve a considerar cualquier PDF modificado en los
últimos tres minutos. La memoria de documentos ya procesados evita que un
archivo abierto se vuelva a abrir.

% ======================================================================
\section{Visor integrado}
% ======================================================================
Los documentos se abren siempre en un visor propio basado en \textbf{Microsoft
WebView2}, el motor de renderizado de PDF de Chromium distribuido con Windows.
El visor presenta las siguientes características:

\begin{itemize}[leftmargin=1.4em,itemsep=2pt]
  \item Carga exclusivamente rutas locales mediante el esquema \texttt{file://};
        no navega a direcciones de Internet.
  \item Detección de cierre exacta a través del evento de cierre de la propia
        ventana.
  \item Apertura rápida: un proceso de WebView2 se mantiene preparado en segundo
        plano para que cada documento se abra sin la demora de un arranque en
        frío.
  \item Cada ventana se ejecuta en su propio hilo, lo que permite mostrar varios
        documentos de forma independiente.
\end{itemize}

% ======================================================================
\section{Filtrado por idioma}
% ======================================================================
El nombre del archivo determina su idioma mediante los sufijos \texttt{\_SPA} y
\texttt{\_ENG}. El comportamiento depende de cuántos documentos se descarguen y
de la preferencia configurada.

\begin{xltabular}{\textwidth}{@{} L L @{}}
\toprule
\textbf{Situación} & \textbf{Comportamiento} \\
\midrule
\endfirsthead
\toprule
\textbf{Situación} & \textbf{Comportamiento} \\
\midrule
\endhead
\bottomrule
\endfoot
Un único documento (con o sin sufijo de idioma) & Se abre siempre. \\
Dos documentos, sin idioma preferido & Se abren ambos. \\
Dos documentos, con idioma preferido & Se abren ambos y, una vez visibles, se cierra automáticamente el que no coincide; permanece el del idioma preferido. \\
\bottomrule
\end{xltabular}

\subsection{Emparejamiento de documentos}
Dos archivos se consideran el mismo documento en distinto idioma cuando, al
retirar el sufijo de idioma y normalizar los separadores, comparten el mismo
nombre base. Por ejemplo, \texttt{D000\_H\_SPA\_Reporte} y
\texttt{D000\_H\_ENG Reporte} se emparejan correctamente.

% ======================================================================
\section{Prioridad por tipo de documento}
% ======================================================================
En producción, un mismo documento puede descargarse en dos formatos: el
\texttt{.pdf} nativo y la versión convertida desde \texttt{.docx}, que el sistema
de origen entrega como \texttt{\_docx.pdf} (la extensión \texttt{.docx} se
convierte en un sufijo \texttt{\_docx}). Pueden llegar así hasta cuatro archivos:
dos idiomas $\times$ dos tipos.

El criterio de prioridad se aplica \textbf{después} del filtro de idioma: para el
mismo documento e idioma, la copia derivada de \texttt{.docx} (\texttt{\_docx.pdf})
prevalece sobre el \texttt{.pdf} nativo, que se cierra automáticamente. La regla es
simétrica, por lo que funciona sin importar cuál de los dos tipos se descargue
primero.

\begin{xltabular}{\textwidth}{@{} L L @{}}
\toprule
\textbf{Situación} & \textbf{Comportamiento} \\
\midrule
\endfirsthead
\toprule
\textbf{Situación} & \textbf{Comportamiento} \\
\midrule
\endhead
\bottomrule
\endfoot
Solo el \texttt{.pdf} nativo & Se abre el nativo. \\
Solo el \texttt{\_docx.pdf} & Se abre el derivado de docx. \\
Ambos tipos del mismo documento e idioma & Permanece el \texttt{\_docx.pdf}; el nativo se cierra. \\
\bottomrule
\end{xltabular}

El sufijo de copia del navegador (\texttt{(1)}, \texttt{(2)}, \ldots) se ignora de
forma uniforme en todas las claves de identidad del archivo. Así, una copia más
reciente reemplaza correctamente a la abierta y el emparejamiento de idioma sigue
funcionando aunque uno de los archivos sea una copia (por ejemplo
\texttt{\_docx (1).pdf}).

% ======================================================================
\section{Reemplazo por copia más reciente}
% ======================================================================
Si un documento ya está abierto y el usuario descarga una copia más reciente del
mismo documento y en el mismo idioma (por ejemplo \texttt{reporte (1).pdf}
mientras \texttt{reporte.pdf} está en pantalla), la aplicación cierra la ventana
de la copia anterior y muestra la nueva. De este modo el operador siempre
visualiza la versión más reciente.

Dado que la copia nueva se abre en una ventana nueva, el límite de visualización
de 20 minutos se reinicia automáticamente. La copia anterior se elimina al
cerrarse, mientras que la copia en uso queda protegida frente a esa eliminación.

% ======================================================================
\section{Límite de tiempo de visualización}
% ======================================================================
Cada documento puede permanecer abierto un máximo de \textbf{20 minutos}. El
control es el siguiente:

\begin{itemize}[leftmargin=1.4em,itemsep=2pt]
  \item A los \textbf{15 minutos} se muestra una notificación que advierte que el
        documento se cerrará en 5 minutos.
  \item A los \textbf{20 minutos} la ventana se cierra automáticamente y el
        documento se elimina como en cualquier cierre.
\end{itemize}

Esta notificación de tiempo es la única notificación de rutina que genera la
aplicación.

% ======================================================================
\section{Eliminación de archivos}
% ======================================================================
La eliminación se limita exclusivamente a la carpeta vigilada y se ejecuta tras
el cierre de cada documento. El proceso contempla los bloqueos temporales
habituales (análisis antivirus, sincronización de carpetas, uso del propio
visor):

\begin{itemize}[leftmargin=1.4em,itemsep=2pt]
  \item \textbf{Reintentos inmediatos:} cada eliminación se reintenta en
        intervalos crecientes durante unos nueve segundos.
  \item \textbf{Limpieza en segundo plano:} si el archivo sigue bloqueado, se
        registra en una lista de pendientes que se reintenta periódicamente hasta
        completarse; ningún archivo queda sin eliminar de forma silenciosa.
  \item \textbf{Protección ante redescarga:} si un archivo pendiente cambia
        (porque se volvió a descargar), se descarta de la lista y se procesa como
        un documento nuevo, evitando borrar contenido vigente.
\end{itemize}

Adicionalmente, la aplicación reconoce las copias duplicadas que generan los
navegadores (por ejemplo, \texttt{documento (2).pdf}) y las incluye en la
limpieza del documento correspondiente.

% ======================================================================
\section{Configuración}
% ======================================================================
La única opción configurable por el usuario es el \textbf{idioma preferido}, que
se selecciona desde la ventana principal de la aplicación. El resto del
comportamiento es fijo por diseño: la carpeta vigilada, el uso del visor
integrado, la eliminación automática y el límite de 20 minutos.

La configuración se guarda en formato JSON, en texto legible:

\begin{lstlisting}[caption={Ejemplo de settings.json.}]
{
  "PreferredLanguage": "SPA"
}
\end{lstlisting}

Si el archivo no existe o está dañado, la aplicación utiliza valores
predeterminados seguros.

\begin{xltabular}{\textwidth}{@{} >{\bfseries}p{3.4cm} p{2.2cm} L @{}}
\toprule
\normalfont\textbf{Parámetro} & \textbf{Valores} & \textbf{Descripción} \\
\midrule
\endfirsthead
\toprule
\normalfont\textbf{Parámetro} & \textbf{Valores} & \textbf{Descripción} \\
\midrule
\endhead
\bottomrule
\endfoot
PreferredLanguage & Any, SPA, ENG & Idioma que se abre cuando se descargan ambas versiones a la vez. \texttt{Any} abre ambos. \\
\bottomrule
\end{xltabular}

% ======================================================================
\section{Almacenamiento de datos}
% ======================================================================
La aplicación almacena sus datos en la carpeta local del usuario, sin afectar a
otros usuarios del equipo.

\begin{xltabular}{\textwidth}{@{} L p{4.2cm} @{}}
\toprule
\textbf{Ruta} & \textbf{Contenido} \\
\midrule
\endfirsthead
\toprule
\textbf{Ruta} & \textbf{Contenido} \\
\midrule
\endhead
\bottomrule
\endfoot
\ruta{\%LOCALAPPDATA\%\bs PdfAutoViewer\bs settings.json} & Configuración del usuario. \\
\ruta{\%LOCALAPPDATA\%\bs PdfAutoViewer\bs WebView2\bs} & Datos de trabajo del visor integrado. \\
\ruta{\%LOCALAPPDATA\%\bs PdfAutoViewer\bs viewer-error.log} & Registro de diagnóstico del visor, si procede. \\
\ruta{\%LOCALAPPDATA\%\bs PdfAutoViewer\bs app-error.log} & Registro de excepciones no controladas, si procede. \\
\bottomrule
\end{xltabular}

% ======================================================================
\section{Requisitos del sistema}
% ======================================================================
\begin{itemize}[leftmargin=1.4em,itemsep=2pt]
  \item Microsoft Windows 10 o 11 (64 bits).
  \item .NET 8 Desktop Runtime.
  \item WebView2 Runtime (incluido de forma predeterminada en Windows 11).
\end{itemize}

% ======================================================================
\section{Compilación y ejecución}
% ======================================================================
El proyecto se compila con el SDK de .NET 8.

\begin{lstlisting}[caption={Compilación y ejecución.}]
# Ejecutar en modo desarrollo
dotnet run --project PdfAutoViewer/PdfAutoViewer.csproj

# Compilar en modo Release
dotnet build -c Release

# Publicar un ejecutable de un solo archivo
dotnet publish PdfAutoViewer/PdfAutoViewer.csproj -c Release ^
       -r win-x64 --self-contained false -o publish
\end{lstlisting}

También puede abrirse la solución \texttt{PdfAutoViewer.sln} en Visual Studio
2022 o JetBrains Rider.

\subsection{Compilación ligera de un solo archivo}
Para distribuir la aplicación como un único ejecutable de aproximadamente
1,8~MB, se empaquetan la app y las DLL de WebView2 (incluido el cargador nativo
\texttt{WebView2Loader.dll}) dentro del \texttt{.exe}, \textbf{sin} incluir .NET.
El equipo destino sigue necesitando el \emph{.NET~8 Desktop Runtime} y el
\emph{WebView2 Runtime}. El script \texttt{publish.bat} ejecuta este mismo
comando.

\begin{lstlisting}[caption={Compilación ligera (un solo archivo, sin .NET embebido).}]
dotnet publish PdfAutoViewer/PdfAutoViewer.csproj -c Release -r win-x64 ^
       --self-contained false ^
       -p:PublishSingleFile=true ^
       -p:IncludeNativeLibrariesForSelfExtract=true ^
       -o publish-light
\end{lstlisting}

El ejecutable se genera en \ruta{publish-light\bs PdfAutoViewer.exe}; los archivos
\texttt{.xml} y \texttt{.pdb} contiguos no se necesitan en ejecución. Para una
compilación totalmente autónoma que \textbf{no} requiera .NET en el destino
(empaqueta el runtime, $\sim$80~MB), use \texttt{-{}-self-contained true}.

\subsection{Instalación del SDK sin privilegios de administrador}
En equipos corporativos sin permisos de administrador, el SDK de .NET 8 puede
instalarse por usuario, en \ruta{\%USERPROFILE\%\bs.dotnet}, y utilizarse con su
ruta completa. Desde la terminal integrada de Visual Studio Code (PowerShell), en
la raíz del proyecto:

\begin{lstlisting}[caption={Ejecución con un SDK local (sin administrador).}]
& "$env:USERPROFILE\.dotnet\dotnet.exe" run --project PdfAutoViewer
& "$env:USERPROFILE\.dotnet\dotnet.exe" test PdfAutoViewer.Tests
\end{lstlisting}

\subsection{Versionado}
La versión de la aplicación se define una sola vez en el archivo de proyecto
(\texttt{PdfAutoViewer.csproj}, etiqueta \texttt{<Version>}) y la ventana de
estado la lee en tiempo de ejecución desde el ensamblado. Para publicar una nueva
versión basta con actualizar ese valor; no es necesario modificar el código.

La aplicación admite una sola instancia por sesión de Windows: si ya hay una
copia en ejecución en la misma sesión, un nuevo lanzamiento se cierra de
inmediato. En entornos de varias sesiones (VDI o multisesión), cada sesión
ejecuta su propia instancia de forma independiente.

% ======================================================================
\section{Dependencias}
% ======================================================================
\begin{xltabular}{\textwidth}{@{} >{\bfseries}p{5.2cm} p{2.4cm} L @{}}
\toprule
\normalfont\textbf{Dependencia} & \textbf{Versión} & \textbf{Función} \\
\midrule
\endfirsthead
\toprule
\normalfont\textbf{Dependencia} & \textbf{Versión} & \textbf{Función} \\
\midrule
\endhead
\bottomrule
\endfoot
.NET (Windows Desktop) & 8.0 & Plataforma de ejecución. \\
Windows Forms & 8.0 & Interfaz de usuario. \\
Microsoft.Web.WebView2 & 1.0.x & Motor del visor integrado. \\
WebView2 Runtime & Sistema & Componente del sistema operativo. \\
\bottomrule
\end{xltabular}

La única dependencia externa de paquetes es \texttt{Microsoft.Web.WebView2};
todas las dependencias son componentes oficiales de Microsoft.

% ======================================================================
\section{Seguridad y privacidad}
% ======================================================================
La aplicación \textbf{no establece conexiones de red}. No incorpora bibliotecas
de cliente HTTP ni mecanismos de telemetría, y el visor integrado carga
únicamente archivos locales mediante el esquema \texttt{file://}. No recopila,
almacena ni transmite información del usuario ni el contenido de los documentos;
de cada archivo solo se consultan su nombre y su tamaño.

La siguiente tabla enumera los recursos del sistema operativo que utiliza la
aplicación, con su nivel de acceso.

\begin{xltabular}{\textwidth}{@{} >{\bfseries}p{4.3cm} p{2.8cm} L @{}}
\toprule
\normalfont\textbf{Recurso} & \textbf{Acceso} & \textbf{Uso} \\
\midrule
\endfirsthead
\toprule
\normalfont\textbf{Recurso} & \textbf{Acceso} & \textbf{Uso} \\
\midrule
\endhead
\bottomrule
\endfoot
Registro de Windows & Lectura & Resolución de la ruta de la carpeta de descargas del usuario. \\
Carpeta vigilada & Lectura / Eliminación & Detección de PDF, lectura de tamaño y eliminación tras el uso. \\
Carpeta de datos de la app & Lectura / Escritura & Configuración y datos de trabajo de WebView2. \\
Procesos & Indirecto & El motor WebView2 ejecuta el renderizado del PDF en procesos propios gestionados por su runtime. \\
Red & Ninguno & La aplicación no realiza conexiones de red. \\
\bottomrule
\end{xltabular}

Se ejecuta con los privilegios del usuario que la inicia; no requiere elevación
de administrador, no instala servicios ni controladores y no modifica
configuraciones del sistema. Sus operaciones de archivo se limitan a la carpeta
vigilada y a su propia carpeta de datos.

% ======================================================================
\section{Manejo de errores y registro}
% ======================================================================
La aplicación está diseñada para degradar de forma controlada:

\begin{itemize}[leftmargin=1.4em,itemsep=2pt]
  \item Si una eliminación falla por un bloqueo temporal, se reintenta en segundo
        plano sin interrumpir la operación.
  \item Los errores relevantes se comunican mediante un aviso en la bandeja del
        sistema, de modo que un fallo no pase inadvertido.
  \item El visor registra los fallos de inicialización en un archivo de
        diagnóstico local para facilitar el soporte.
  \item Las excepciones no controladas se registran en un archivo de diagnóstico
        local; las del hilo de interfaz se capturan para que un error puntual no
        detenga la aplicación.
\end{itemize}

% ======================================================================
\section{Limitaciones conocidas}
% ======================================================================
\begin{itemize}[leftmargin=1.4em,itemsep=3pt]
  \item El filtrado por idioma solo se activa cuando ambas versiones se descargan
        de forma prácticamente simultánea; descargas separadas en el tiempo se
        tratan como documentos independientes.
  \item El reemplazo por copia más reciente aplica a documentos del mismo idioma;
        una versión en otro idioma se gestiona mediante el filtro de idioma.
\end{itemize}

% ======================================================================
\section{Historial de versiones}
% ======================================================================
\begin{xltabular}{\textwidth}{@{} p{1.8cm} p{2.6cm} L @{}}
\toprule
\textbf{Versión} & \textbf{Fecha} & \textbf{Cambios} \\
\midrule
\endfirsthead
\toprule
\textbf{Versión} & \textbf{Fecha} & \textbf{Cambios} \\
\midrule
\endhead
\bottomrule
\endfoot
2.2.0 & Junio de 2026 & Prioridad por tipo: el PDF derivado de \texttt{.docx} (\texttt{\_docx.pdf}) prevalece sobre el \texttt{.pdf} nativo tras el filtro de idioma; manejo uniforme del sufijo de copia \texttt{(n)} en todas las claves de identidad; \texttt{publish.bat} genera la compilación ligera de un solo archivo. \\
2.1.0 & Junio de 2026 & Identidad de marca «Philips Document Flow (PDF)»; icono azul con el texto «PDF» en el ejecutable, la ventana y la bandeja; interfaz totalmente en inglés; versión leída del ensamblado; manual técnico en español e inglés. \\
2.0 & --- & Visor integrado único (WebView2); configuración mínima; límite de visualización de 20 minutos; robustez para operación continua. \\
\bottomrule
\end{xltabular}

\end{document}
